Put \(D=y+E\).  First we construct the exceptional set.  Apply \(\text{lem:semialgebraic-stratification-sard}\) to choose finite Whitney stratifications of \(C\) and \(\partial E\) by Nash manifolds, denoted by \((S_j)_j\) and \((R_\ell)_\ell\).  For every pair of strata define
\[
 F_{j\ell}:S_j\times R_\ell\longrightarrow\mathbb R^d,
 \qquad F_{j\ell}(c,e)=c-e.
\]
The product \(S_j\times R_\ell\) is a Nash manifold, and \(F_{j\ell}\) is Nash because it is the restriction of the polynomial subtraction map.  Hence the narrowed Nash Sard assertion in \(\text{lem:semialgebraic-stratification-sard}\) makes the set \(B_{j\ell}\) of critical values of \(F_{j\ell}\) semialgebraic of dimension strictly smaller than \(d\).  Such a set has empty interior and is contained in its boundary; \(\text{lem:semialgebraic-boundary-null}\) therefore gives it Lebesgue measure zero.  The finite union
\[
 B=\bigcup_{j,\ell}B_{j\ell}
\]
is semialgebraic and Lebesgue null.

Fix \(y\notin B\).  If \(x\in C\cap(y+\partial E)\), let \(S_j\) and \(R_\ell\) contain \(x\) and \(x-y\), respectively.  Since \(y=F_{j\ell}(x,x-y)\) is a regular value,
\[
 T_xS_j+T_{x-y}R_\ell=\mathbb R^d. \tag{1}
\]
We spell out the consequence of (1) that is used below.  In Nash charts around these two points, choose a linear complement of \(T_{x-y}R_\ell\) contained in the image of \(T_xS_j\).  The derivative of the corresponding normal-coordinate projection from \(S_j\) is onto by (1).  The local submersion normal form therefore makes the slice \(S_j-y\) cross every sufficiently small normal slice to \(R_\ell\).  Because the chosen stratification is Whitney and \(E=\overline{\operatorname{int}E}\), the local normal slices of \(E\) are the closures of their relative interiors.  Thus there are \(x_r\in S_j\cap(y+\operatorname{int}E)\) with \(x_r\to x\).  This argument also covers a zero-dimensional boundary stratum: then (1) forces \(T_xS_j=\mathbb R^d\), so \(S_j\) is locally open.  Points of \(C\cap(y+\operatorname{int}E)\) already belong to the set being closed.  Consequently
\[
 C\cap D=\overline{C\cap\operatorname{int}D}. \tag{2}
\]

We next prove the intersection limit under the stated upper and erosion hypotheses.  Let
\(A_n=\overline{C_n\cap D_n}\) and \(A=C\cap D\).  For the upper limit, take a subsequence and \(x_n\in A_n\) with \(x_n\to x\).  Choose \(z_n\in C_n\cap D_n\) with \(\lVert z_n-x_n\rVert\le n^{-1}\).  Local Painleve--Kuratowski upper convergence of \(C_n\), together with \(\limsup_nD_n\subseteq D\), gives
\[
 x=\lim_nz_n\in C\cap D=A. \tag{3}
\]

For the lower limit, fix \(x\in A\).  By (2), choose \(x^{(r)}\in C\cap\operatorname{int}D\) such that \(x^{(r)}\to x\).  For each \(r\), take \(\varepsilon_r>0\) with
\(K_r=\overline{B}(x^{(r)},\varepsilon_r)\subset\operatorname{int}D\).  The local Painleve--Kuratowski lower convergence supplies \(x_n^{(r)}\in C_n\) with \(x_n^{(r)}\to x^{(r)}\).  Eventually \(x_n^{(r)}\in K_r\), while the erosion hypothesis gives \(K_r\subseteq D_n\) eventually.  Hence \(x_n^{(r)}\in C_n\cap D_n\) eventually.  A diagonal choice in \(n\) and then \(r\) yields points of \(A_n\) converging to \(x\), and therefore
\[
 A\subseteq\liminf_n A_n. \tag{4}
\]
Equations (3)--(4) give Painleve--Kuratowski, hence Fell, convergence.  By the common-compact-bound hypothesis, all sufficiently large \(A_n\) lie in one compact set.  If \(A=\varnothing\), (3) and compactness force \(A_n=\varnothing\) eventually.  If \(A\ne\varnothing\), (3)--(4) on that compact ambient space are equivalent to Hausdorff convergence.  This proves every assertion of the corrected statement.